\documentclass[11pt,a4paper]{amsart}
\usepackage[margin=30mm]{geometry}
\usepackage{amsmath,amssymb,amsthm,mathtools}
\usepackage{enumitem}
\usepackage{booktabs,longtable}
\usepackage[hidelinks]{hyperref}
\usepackage{url}

\newcommand{\Q}{\mathbb{Q}}
\newcommand{\Z}{\mathbb{Z}}
\newcommand{\F}{\mathbb{F}}
\newcommand{\Qp}{\mathbb{Q}_p}
\newcommand{\Zp}{\mathbb{Z}_p}
\newcommand{\Fp}{\mathbb{F}_p}
\newcommand{\Reg}{\mathrm{Reg}}
\newcommand{\tors}{\mathrm{tors}}
\newcommand{\vp}{v_p}
\newcommand{\qp}{q_p}
\newcommand{\Ebar}{\widetilde{E}}
\newcommand{\Lam}{\Lambda}
\makeatletter\g@addto@macro\UrlBreaks{\do\_\do\a\do\e\do\i\do\o\do\u}\makeatother
\newcommand{\lean}[1]{\nolinkurl{#1}}

\theoremstyle{plain}
\newtheorem*{mainthm}{Main Theorem}
\newtheorem{thmx}{Theorem}
\renewcommand{\thethmx}{\Alph{thmx}}
\newtheorem{lemma}{Lemma}
\newtheorem{proposition}[lemma]{Proposition}
\newtheorem*{thmB}{Theorem B}
\newtheorem*{thmC}{Theorem C}
\newtheorem*{corD}{Corollary D}
\newtheorem*{corE}{Corollary E}
\newtheorem*{thmF}{Theorem F}
\newtheorem*{corFp}{Corollary F$'$}
\newtheorem*{thmH}{Theorem H}
\newtheorem*{corHp}{Corollary H$'$}
\newtheorem*{thmJ}{Theorem J}
\newtheorem*{thmS}{Theorem S}
\newtheorem*{thmR}{Theorem R}
\newtheorem*{propRpp}{Proposition R$''$}
\newtheorem*{corRp}{Corollary R$'$}
\newtheorem*{thmRs}{Theorem R$_s$}
\newtheorem*{corRps}{Corollary R$'_s$}
\newtheorem*{propRpps}{Proposition R$''_s$}
\newtheorem*{propRplus}{Proposition R$^+$}
\newtheorem*{prop6}{Proposition 6}
\newtheorem*{corG}{Corollary G}
\theoremstyle{remark}
\newtheorem*{remark}{Remark}

\begin{document}

\title[Integrality of the $p$-adic height at anomalous primes]{Integrality of the $p$-adic height pairing at an anomalous prime\\ and points of order $p$ in $E(\Qp)$}
\author{Yoshikazu Hosomura}
\address{Independent researcher, Japan}
\email{wvbks0@gmail.com}
\date{\today}
\subjclass[2020]{11G50, 11G07, 11S31 (primary); 11Y40, 68V20 (secondary)}
\keywords{$p$-adic height, $p$-adic regulator, anomalous prime, local torsion, division polynomial, Fermat quotient, Lean 4}

\begin{abstract}
Let $E/\Q$ be an elliptic curve and $p\ge 5$ a prime of good ordinary reduction. The cyclotomic $p$-adic regulator $\Reg_p$ is in general not $p$-integral when $p$ is \emph{anomalous}, $\#\Ebar(\Fp)=p$; counting indices gives the elementary bound $\vp(\Reg_p)\ge r-2$, and computations of $p$-adic heights have usually excluded such primes. We show that the deficit is governed by a single local condition. If $p\nmid\prod_\ell c_\ell$ and $E(\Q)/\tors$ surjects onto $\Ebar(\Fp)$, then \emph{every entry of the height matrix is $p$-integral if and only if $E(\Qp)$ contains a point of order $p$}; in rank one $\vp(\Reg_p)\ge0$ exactly in that case, and otherwise $\vp(\Reg_p)=-1$. Modulo $p$ the height matrix has rank exactly one in the non-split case and rank at most two in the split case, its constant being a Fermat quotient, respectively a second-order Fermat quotient, of the $p$-th division polynomial at one point. Since points of order $p$ in $E(\Qp)$ are rare by a theorem of David and Weston, the non-split behaviour is the generic one at anomalous primes of non-CM curves, while CM curves are split at every anomalous prime; in particular the result applies at the anomalous primes excluded in recent work on rank-two CM curves. The arithmetic steps of the proofs have been formalised in Lean~4.
\end{abstract}

\maketitle

\section{Introduction}\label{sec:intro}

\subsection{The question}
Let $E/\Q$ be an elliptic curve given by a global minimal Weierstrass equation, $p\ge5$ a prime of good ordinary reduction, and $\Lam:=E(\Q)/\tors$, of rank $r$. The cyclotomic $p$-adic height of Mazur--Tate and Schneider \cite{MT91,Sch85} is a quadratic form $h_p$ on $\Lam$ with values in $\Qp$; its discriminant $\Reg_p$ on a basis of $\Lam$ is the regulator in the $p$-adic analogue of the conjecture of Birch and Swinnerton-Dyer \cite{MTT,BMS}. Unlike the archimedean regulator, $\Reg_p$ is not known to be non-zero (Schneider's conjecture), and we say nothing about that question. This paper concerns a coarser invariant: the valuation $\vp(\Reg_p)$, and more precisely the $p$-integrality of the individual entries of the height matrix.

When $p$ is \emph{anomalous} in the sense of Mazur \cite{Maz72}, i.e.\ $p\mid\#\Ebar(\Fp)$ (for $p\ge7$ this means $a_p=1$, $\#\Ebar(\Fp)=p$), the regulator need not be $p$-integral. Negative valuations are visible in the tables of Mazur, Stein and Tate \cite[\S\S4.2--4.6]{MST}, where the anomalous primes are noted, and the loss of precision they cause is familiar to anyone computing $p$-adic heights \cite{Har}. A simple count of indices explains the size of the deficit: $\vp(\Reg_p)\ge r-2$ (Theorem~\ref{thm:A} below), a bound which is elementary and which we do not claim. Work that relies on integrality therefore excludes anomalous primes. Balakrishnan, \c{C}iperiani and Stein restrict to ordinary non-anomalous primes not dividing the Tamagawa product, where the $p$-adic height pairing is integral \cite{BCS}; examples at anomalous primes are discussed separately by Stein and Wuthrich \cite{SW}; and in the recent work of Banwait \cite{Ban} on the Shafarevich--Tate group of rank-two CM elliptic curves, anomalous primes form one of the classes excluded throughout, non-anomalousness being what extends a regulator estimate from $E^0(\Q)$ to $E(\Q)$ \cite[\S2.2]{Ban}.

The purpose of this paper is to describe exactly what happens at the excluded primes, and thereby to complement such work.

\subsection{Results}
Call $E$ \emph{split at $p$} when $E(\Qp)[p]\ne0$; for an anomalous prime this says that the reduction sequence $0\to E_1(\Qp)\to E(\Qp)\to\Ebar(\Fp)\to0$ splits. Write $c_\ell$ for the Tamagawa numbers and let $\iota$ denote the hypothesis that the reduction map $\Lam\to\Ebar(\Fp)$ is onto.

\begin{mainthm}[Theorem~\ref{thm:main}]
Suppose $\#\Ebar(\Fp)=p$, $p\nmid\prod_\ell c_\ell$ and $\iota$. Let $M_\Lam$ be the height matrix on a basis of $\Lam$.
\begin{enumerate}[label=\textup{(\roman*)}]
\item Every entry of $M_\Lam$ lies in $\Zp$ if and only if $E(\Qp)[p]\ne0$.
\item If $r=1$, then $\vp(\Reg_p)\ge0$ if and only if $E(\Qp)[p]\ne0$; in the non-split case $\vp(\Reg_p)=-1$ exactly.
\item If $E(\Qp)[p]\ne0$, then $\vp(\Reg_p)\ge\max(r-2,0)$ in every rank.
\end{enumerate}
\end{mainthm}

The statement about the \emph{entries} in (i) is not implied by any bound on the determinant and holds in every rank. The bound on the determinant improves on the elementary one in rank one only, where (ii) says exactly which curves have a non-integral regulator.

The same mechanism determines the whole height matrix modulo $p$ (Section~\ref{sec:modp}). In the non-split case $p\,M_\Lam$ is integral and $p\,M_\Lam\bmod p=c\,\chi\chi^{t}$ is of rank exactly one, $\chi$ being the reduction character and $c=2\lambda_1$ a Fermat quotient of the $p$-th division polynomial at one point (Theorem~R, Proposition~R$''$). In the split case $M_\Lam\bmod p$ vanishes on $\ker\chi$ and has rank at most two, and its diagonal constant is a second-order Fermat quotient of the same polynomial (Theorem~R$_s$, Proposition~R$''_s$). In both cases $\vp(\Reg_p)=(r-2)+\vp(\det T)$ for an integral Schur complement $T$ (Corollaries~R$'$, R$'_s$).

\subsection{The mechanism}
The proof is local and runs through one element of $\Fp$. For a point $P'\in E(\Q)$ whose reduction has exact order $p$, let $\phi_p$ be the $p$-th division polynomial and put
\[
\lambda:=\tfrac12\,\qp\bigl(\phi_p(x(P'))\bigr)\in\Fp,\qquad \qp(u):=\frac{u^{p-1}-1}{p}\bmod p .
\]
Writing the height of $pP'$ as a formal-group part plus $2\log_p$ of a ratio of coordinates (Lemma~\ref{lem:decomp}), the second $p$-adic digit of the height is $-\lambda$ (Theorem~B). In the split case $\lambda$ vanishes because Vélu's factorisation of multiplication by $p$ through a rational subgroup of order $p$ turns $\phi_p(x(T))$ into a $p$-th power (Theorem~C, Corollary~D). Conversely $\lambda=0$ forces splitting (Proposition~6). The latter equivalence is a known criterion for local torsion: it is the case $n=1$ of Gross's tameness criterion \cite{Gro90} (as quoted in \cite[Thm.~4.1]{LR}), and it is contained in the work of David and Weston on local torsion primes \cite[Lemma~3.2]{DW} and of Ozeki and Yoshida \cite[Lemma~4.1]{OY}, in the form ``$E$ is split at $p$ if and only if $E\bmod p^2$ is the canonical lift of $\Ebar$''. What we add is its presentation as the vanishing of a Fermat quotient of $\phi_p$ at a single point, which is what connects it to the height, together with an independent proof (Appendix~\ref{app:prop6}).

\subsection{How often each case occurs}
By a theorem of David and Weston, an elliptic curve over $\Q$ has, on average, only finitely many primes $p$ with $E(\Qp)[p]\ne0$ \cite{DW}; for non-CM curves the anomalous primes themselves have density zero. Thus, for a non-CM curve, the anomalous primes at which the Main Theorem applies are, apart from finitely many on average, in the non-split case, where in rank one $\vp(\Reg_p)=-1$ exactly. For CM curves the opposite holds: a CM curve is split at every anomalous prime of good ordinary reduction (\cite[Prop.~2.1(3)]{DW}; Theorem~F below), so a CM curve never shows the deficit (Corollary~F$'$). In particular, for the rank-two CM curves studied in \cite{Ban}, at an anomalous prime satisfying the hypotheses above the regulator is integral, and Corollary~R$'_s$ says that it is a unit exactly when the height matrix modulo $p$ has rank two (Corollary~G in Section~\ref{sec:CM}). We do not extend the criterion of \cite{Ban} to anomalous primes, which would require in addition the anomalous Euler factor of the $p$-adic $L$-function; we only describe the regulator there.

\subsection{Formalisation and computations}
The arithmetic steps of the proofs --- the Fermat-quotient bookkeeping, the determinant estimates, the bilinear algebra behind Theorems~R and R$_s$, and the weight and residue counts of Appendix~\ref{app:prop6} --- have been formalised in Lean~4 with Mathlib, with no \texttt{sorry}, no \texttt{native\_decide}, and axioms among \texttt{propext}, \texttt{Classical.choice}, \texttt{Quot.sound}. The geometric inputs (the $p$-adic sigma function, division polynomials, Vélu's formulas, formal groups, the Tate curve) are not in Mathlib and enter the Lean statements as hypotheses; Appendix~\ref{app:lean} lists the correspondence. The statements were also checked numerically with PARI/GP on curves from Cremona's tables; these checks are reported only where they illustrate a statement.

\subsection*{Organisation}
Section~\ref{sec:setting} fixes the normalisation. Section~\ref{sec:baseline} gives the elementary baseline and the decomposition of the height. Section~\ref{sec:criterion} proves the criterion through $\lambda$, Section~\ref{sec:main} the Main Theorem, Section~\ref{sec:modp} the structure modulo $p$, and Section~\ref{sec:CM} the CM case and the relation with \cite{Ban}. Section~\ref{sec:related} discusses related work. Appendix~\ref{app:prop6} contains the proof of Proposition~6 and Appendix~\ref{app:lean} the Lean correspondence.

\subsection*{Relation to an earlier version}
This paper is a revised and condensed version of the first part of a longer draft that the author published under the pen name \emph{Kiichi} on the website \emph{computo ergo sum} and deposited on Zenodo \cite{Draft}. The remaining parts of that draft (the higher $p$-adic digits of the height and an archimedean threshold) are treated in a companion paper \cite{Companion}.

\section{Setting and normalisation}\label{sec:setting}

Throughout, $E/\Q$ is given by a global minimal Weierstrass equation with coefficients $a_1,\dots,a_6$, and $p\ge5$ is a prime of good ordinary reduction. For $P\in E(\Q)$, $P\ne O$, we write
\[
P=\bigl(\alpha(P)/d(P)^2,\ \beta(P)/d(P)^3\bigr),\qquad \gcd(\alpha,d)=\gcd(\beta,d)=1,\ d\ge1 .
\]
Let $t=-x/y$ be the parameter of the formal group $\widehat E$ and $E_n(\Qp):=\{P:\vp(t(P))\ge n\}$. The valuation $\vp$ is normalised by $\vp(p)=1$; $\log_p$ is the Iwasawa logarithm with $\log_p(p)=0$. For $u\in\Zp^\times$ the Fermat quotient $\qp(u)$ is defined as above; $\qp\colon\Zp^\times\to\Fp$ is a surjective homomorphism and
\begin{equation}\label{eq:logq}
\vp(\log_p u)=1+\vp(\qp(u)),\qquad \log_p u\equiv -p\,\qp(u)\pmod{p^2}.
\end{equation}
So ``the second $p$-adic digit of $\log_p u$'' and ``the Fermat quotient of $u$'' are the same datum.

\subsection*{The height}
Let $\sigma_p\in t\Zp[[t]]$ be the Mazur--Tate $p$-adic sigma function \cite{MT91}, the unique odd series $\sigma(t)=t+c_2t^2+\cdots$ with $p$-integral coefficients solving the differential equation of \cite{MT91}; here $c_2=a_1/2$. Following \cite[\S2.2]{Har}, if $P\in E(\Q)$ satisfies
\begin{itemize}
\item[(A1)] $P$ reduces to $O$ at $p$, and
\item[(A2)] $P$ reduces to a non-singular point at every bad prime,
\end{itemize}
we put $h_p(P):=2\log_p\bigl(\sigma_p(P)/d(P)\bigr)$, and extend to all of $E(\Q)$ by $h_p(Q):=h_p(mQ)/m^2$ for any $m\ge1$ with $mQ$ satisfying (A1) and (A2). This is $2p$ times the normalisation of \cite{MST}, so the two valuations of $\Reg_p$ differ by $r$; it is the normalisation used in \cite{Har}. The associated pairing is $h_p(P,R):=\tfrac12\bigl(h_p(P+R)-h_p(P)-h_p(R)\bigr)$ and $\Reg_p:=\det\bigl(h_p(P_i,P_j)\bigr)$ for a basis $(P_i)$ of $\Lam$. Valuations of heights and of $\Reg_p$ do not depend on the sign convention; with this normalisation $h_p$ is the negative of PARI's canonical $p$-adic height.

\subsection*{Anomalous, split, $\iota$}
The prime $p$ is \emph{anomalous} if $p\mid\#\Ebar(\Fp)$. By Hasse's bound $\#\Ebar(\Fp)\in\{p,2p\}$ in that case, and $2p$ occurs only for $p=5$ (see Proposition~R$^+$). From Section~\ref{sec:criterion} on we assume $\#\Ebar(\Fp)=p$. The curve is \emph{split at $p$} if $E(\Qp)[p]\ne0$. The hypothesis $\iota$ says that $\Lam\to\Ebar(\Fp)$ is onto. We write $\Lam_0:=\ker(\Lam\to\Ebar(\Fp))$ and $n:=\prod_\ell c_\ell$.

\subsection*{Division polynomials}
$\phi_m,\psi_m,\omega_m$ are normalised by $x(mQ)=\phi_m/\psi_m^2$, $y(mQ)=\omega_m/\psi_m^3$, with $\phi_m$ monic of degree $m^2$ in $x$. Modulo $p$ and for $E$ ordinary, $\phi_p(x)\equiv f(x^p)$ for a polynomial $f$, because $[p]$ factors through Frobenius in characteristic $p$.

\subsection*{The model}
The height and its valuation are invariants of $E/\Q$, but the two parts of the decomposition in Lemma~\ref{lem:decomp} depend on the Weierstrass model. The substitution $(x,y)\mapsto(x,y+sx)$ with $s\in\Z$ carries a global minimal equation to a global minimal one, fixes $x$, $\alpha$, $d$, $\phi_p$, $\psi_p$ and $\lambda$, and changes $a_1$ to $a_1+2s$. Since $p$ is odd we may and do choose $p\mid a_1$ whenever it is convenient; statements about the individual parts refer to such a model.

\section{The baseline and the decomposition of the height}\label{sec:baseline}

\begin{thmx}[the baseline]\label{thm:A}
Let $G:=\{P\in E(\Q): P \text{ satisfies (A1) and (A2)}\}$ and $k:=[\Lam:\mathrm{image}\ \mathrm{of}\ G]$. Then $\vp(\Reg_p)\ge r-2\vp(k)$. If moreover $p\nmid n$, then $\vp(k)\le1$, with equality exactly when the image of $\Lam$ in $\Ebar(\Fp)$ meets $\Ebar(\Fp)[p]\ne0$. In particular $\vp(\Reg_p)\ge r-2$ in that case and $\vp(\Reg_p)\ge r$ otherwise.
\end{thmx}

This bound is elementary and known to practitioners (see Section~\ref{sec:related}); we include the proof because the rest of the paper is measured against it.

\begin{lemma}\label{lem:unit}
If $P$ satisfies \textup{(A1)} and \textup{(A2)} and $e:=\vp(t(P))\ge1$, then $\sigma_p(P)/d(P)\in\Zp^\times$, $h_p(P)\in 2p\Zp$, and $\vp(h_p(P))=1+\vp\bigl(\qp(\sigma_p(P)/d(P))\bigr)$.
\end{lemma}

\begin{proof}
Since $\sigma_p(t)/t\in1+t\Zp[[t]]$, $\vp(\sigma_p(P))=e$. From $\vp(x(P))=-2e$ and $\gcd(\alpha,d)=1$ we get $\vp(d)=e$ and $\vp(\alpha)=0$, hence $\sigma_p(P)/d(P)\in\Zp^\times$. The remaining assertions are \eqref{eq:logq}.
\end{proof}

\begin{lemma}[decomposition]\label{lem:decomp}
Let $P$ satisfy \textup{(A1)}, \textup{(A2)} and $e=\vp(t(P))$. Put $u(P):=\sigma_p(t)/t\in1+p^e\Zp$ and $n_\ell(P):=v_\ell(\alpha)-v_\ell(\beta)\in\Z$ for primes $\ell\ne p$. Then
\[
h_p(P)=2\Lambda_\sigma+2A,\qquad \Lambda_\sigma:=\log_p u(P),\quad A:=\log_p\bigl(\alpha(P)/\beta(P)\bigr)=\sum_{\ell\ne p}n_\ell(P)\log_p\ell .
\]
With $b:=\vp(\Lambda_\sigma)$ and $c:=\vp(A)$ one has $c\ge1$, $b\ge e$, and $b\ge\min(e+\vp(a_1),\,2e)$. In particular, in a model with $p\mid a_1$, $b\ge2$ when $e=1$ and $b\ge3$ when $e\ge2$.
\end{lemma}

\begin{proof}
From $t=-x/y$, $x=\alpha/d^2$, $y=\beta/d^3$ we get $t/d=-\alpha/\beta$, hence $\sigma_p(P)/d(P)=u(P)\cdot(-\alpha/\beta)$. Taking $\log_p$ and using $\log_p(-1)=0$ gives the decomposition; $\alpha$ and $\beta$ are prime to $p$ by Lemma~\ref{lem:unit}, so $A$ is a $\Z$-combination of $\log_p\ell$ with $\ell$ a $p$-adic unit, and $c\ge1$ by \eqref{eq:logq}. Finally $u=1+c_2t+O(t^2)$ with $c_2=a_1/2$ and $p$-integral higher coefficients, so $\vp(u-1)\ge\min(e+\vp(a_1),2e)$, and $\vp(\log_p u)=\vp(u-1)$ since $p\ge3$.
\end{proof}

\begin{proof}[Proof of Theorem~\ref{thm:A}]
If $P,R,P+R\in G$, Lemma~\ref{lem:unit} puts the three heights in $2p\Zp$, so $h_p(P,R)\in p\Zp$ as $p$ is odd. The height matrix on a basis of the image of $G$ therefore has entries in $p\Zp$ and its determinant has valuation $\ge r$. A change of basis from $\Lam$ to the image of $G$ has determinant $\pm k$, so this determinant is $k^2\Reg_p$, which gives the first bound. The quotient of $\Lam$ by the image of $G$ embeds into $\Ebar(\Fp)\times\prod_\ell\Phi_\ell(\F_\ell)$; if $p\nmid n$ its $p$-part comes from $\Ebar(\Fp)$, whose $p$-part is cyclic of order at most $p$ by Hasse's bound.
\end{proof}

\section{The criterion: one element of \texorpdfstring{$\Fp$}{Fp}}\label{sec:criterion}

From now on $\#\Ebar(\Fp)=p$. Let $P'\in E(\Q)$ have reduction of exact order $p$ and satisfy (A2), put $P:=pP'$ and $e:=\vp(d(P))=\vp(t(P))\ge1$.

\begin{lemma}[splitting and depth]\label{lem:depth}
$E(\Qp)[p]\ne0$ if and only if $e\ge2$.
\end{lemma}

\begin{proof}
Suppose $\tau\in E(\Qp)$ has order $p$. For $p\ge5$ the formal group $\widehat E(p\Zp)$ is torsion-free, so $\tau$ does not reduce to $O$, and since $\#\Ebar(\Fp)=p$ its reduction generates $\Ebar(\Fp)$. Hence $R:=P'-j\tau\in E_1(\Qp)$ for some $j$, and $pP'=pR$. As $[p](t)=pt+(\text{terms of degree}\ge2)$ with $\Zp$-coefficients, $\vp([p](t))\ge\min(1+\vp(t),2\vp(t))\ge2$ for $\vp(t)\ge1$, so $e\ge2$.

Conversely suppose $E(\Qp)[p]=0$. Then $E(\Qp)$ is torsion-free (its prime-to-$p$ torsion injects into $\Ebar(\Fp)$, of order $p$), so it is a torsion-free compact $p$-adic Lie group of dimension one with $E_1(\Qp)$ of index $p$; hence $E(\Qp)\cong\Zp$ and $E_1(\Qp)=pE(\Qp)$. The formal logarithm $\log_E\colon E_1(\Qp)\to p\Zp$ is an isomorphism for $p\ge5$, and $\ell(Q):=\log_E(pQ)/p$ extends it to an injective homomorphism $E(\Qp)\to\Zp$ with $\ell(E(\Qp))=\Zp$ and $\ell(E_1(\Qp))=p\Zp$. As $P'\notin E_1(\Qp)$, $\ell(P')\in\Zp^\times$, so $\vp(\log_E(pP'))=1$ and $e=1$.
\end{proof}

\begin{lemma}[$\lambda$ is local]\label{lem:local}
$\phi_p$ modulo $p$ is a polynomial in $x^p$, so $\phi_p(x)\bmod p^2$ depends only on $x\bmod p$. Consequently
\[
\lambda:=\tfrac12\,\qp\bigl(\phi_p(x(P'))\bigr)\in\Fp
\]
is a function $\lambda(\bar x)$ of the reduction of $E$ modulo $p^2$ and of the class $\bar x$ of $x(P')$ in $\Ebar(\Fp)\smallsetminus\{O\}$ alone.
\end{lemma}

\begin{proof}
In characteristic $p$, $[p]=V\circ F$, so $x\circ[p]$ is a rational function of $x^p$ and $\phi_p'\equiv0\pmod p$. For a polynomial $g$ with $p\mid g'(x)$, Taylor's formula gives $p^2\mid g(x+pk)-g(x)$ for $x,k\in\Zp$. That $\phi_p(x(P'))$ is a unit is shown in the proof of Theorem~B.
\end{proof}

\begin{thmB}[criterion]
If $e\ge2$, or if the model satisfies $p\mid a_1$, then $\qp\bigl(\alpha(P)/\beta(P)\bigr)=-\lambda$. Hence, in a model with $p\mid a_1$, $c\ge2$ if and only if $\lambda=0$.
\end{thmB}

\begin{proof}
\emph{Step 1 (denominators).} Write $x(P')=a/d'^2$, $y(P')=b/d'^3$ with $p\nmid d'$, possible since $P'$ does not reduce to $O$. The homogenised values $\Phi:=d'^{2p^2}\phi_p(P')$, $\Psi:=d'^{p^2-1}\psi_p(P')$ and $\Omega:=d'^{3p^2}\omega_p(P')$ are integers, and $x(P)=\Phi/(d'\Psi)^2$, $y(P)=\Omega/(d'\Psi)^3$. By (A2) there is no cancellation, so $d(P)=d'|\Psi|$, $\alpha=\pm\Phi$, $\beta=\pm\Omega$ and $\alpha/\beta=(\phi_p(P')/\omega_p(P'))\,d'^{-p^2}$. As $\qp$ is a homomorphism and $\qp(d'^{p^2})=0$,
\begin{equation}\label{eq:star}
\qp(\alpha/\beta)=\qp\bigl(\phi_p(P')\bigr)-\qp\bigl(\omega_p(P')\bigr).
\end{equation}
In particular $\phi_p(P')$ and $\omega_p(P')$ are $p$-adic units.

\emph{Step 2 (the Weierstrass equation).} Write $\phi=\phi_p(P')$, $\omega=\omega_p(P')$, $\psi=\psi_p(P')$. Multiplying the equation at $pP'$ by $\psi^6$,
\[
\omega^2+a_1\phi\omega\psi+a_3\omega\psi^3=\phi^3+a_2\phi^2\psi^2+a_4\phi\psi^4+a_6\psi^6 .
\]
Since $\vp(\psi)=\vp(d(P))=e$, every term of $\omega^2-\phi^3$ other than $a_1\phi\omega\psi$ has valuation $\ge2e\ge2$, and that one has valuation $\vp(a_1)+e$. Under either hypothesis $\omega^2\equiv\phi^3\pmod{p^2}$.

\emph{Step 3.} Write $\phi^{p-1}=1+pA'$ and $\omega^{p-1}=1+pB'$. Raising $\omega^2\equiv\phi^3\pmod{p^2}$ to the power $p-1$ gives $2B'\equiv3A'\pmod p$, so $\qp(\phi)=2\lambda$ and $\qp(\omega)=3\lambda$, and \eqref{eq:star} gives $\qp(\alpha/\beta)=-\lambda$. The last assertion follows from \eqref{eq:logq}: $c\ge2$ if and only if $\qp(\alpha/\beta)=0$.
\end{proof}

\begin{remark}
Without either hypothesis the term $a_1\phi\omega\psi$ survives modulo $p^2$ and one finds $\qp(\alpha/\beta)=-\lambda-w_1/2$ with $w_1:=a_1\,\omega\,(\psi/p)/\phi^2\bmod p$. This is why the model with $p\mid a_1$ is used.
\end{remark}

\begin{thmC}[Vélu factorisation]
Let $K$ be a field of characteristic $0$, $E/K$ given by a Weierstrass equation, $m\ge3$ odd, $D\subset E(\bar K)$ a Galois-stable cyclic subgroup of order $m$, and $T\in D\smallsetminus\{O\}$ with $x(T)\in K$. Let
\[
\Pi(x):=\prod_{\{\pm Q\}\subset D\smallsetminus\{O\}}\bigl(x-x(Q)\bigr),
\]
monic of degree $(m-1)/2$, and let $\bar\phi$, monic of degree $m$, be the numerator of the $x$-part $\bar\phi/\Pi^2$ of the Vélu isogeny $\pi_D\colon E\to E/D$. Then $\phi_m(x(T))=\bar\phi(x(T))^m$.
\end{thmC}

\begin{proof}
$D$ being Galois-stable, $\pi_D$ and its dual $\hat\pi_D$ are defined over $K$ and $[m]=\hat\pi_D\circ\pi_D$. Fix a model of $E/D$ and write $x\circ\hat\pi_D=\phi^\vee/(\psi^\vee)^2$ with $\phi^\vee(X)=\sum_{i\le m}c_iX^i$, $c_m=:A\ne0$, $\deg(\psi^\vee)^2=m-1$. Clearing denominators in $x\circ[m]=(\phi^\vee/(\psi^\vee)^2)\circ(\bar\phi/\Pi^2)$ gives $\phi_m/\psi_m^2=N/\mathrm{Den}$ with
\[
N(x):=\sum_{i\le m}c_i\,\bar\phi(x)^i\,\Pi(x)^{2(m-i)},\qquad \mathrm{Den}(x):=\Pi^{2m}\,\psi^\vee(\bar\phi/\Pi^2)^2 .
\]
Since $\deg(\bar\phi^i\Pi^{2(m-i)})=m^2-m+i$, only $i=m$ reaches the top degree: $\deg N=m^2$ with leading coefficient $A$, and $\deg\mathrm{Den}=m^2-1$. As $x\circ[m]$ has degree $m^2$, the fraction $N/\mathrm{Den}$ is reduced, and since $\phi_m$, $\psi_m^2$ are coprime with $\phi_m$ monic, $N=A\phi_m$. At $x=x(T)$ one has $\Pi(x(T))=0$, so $N(x(T))=A\,\bar\phi(x(T))^m$, and $A$ cancels.
\end{proof}

\begin{corD}
If $E(\Qp)[p]\ne0$, then $\lambda(\bar x)=0$ for every $\bar x\in\Ebar(\Fp)\smallsetminus\{O\}$.
\end{corD}

\begin{proof}
Let $\tau\in E(\Qp)$ have order $p$; as in Lemma~\ref{lem:depth}, $\tau$ does not reduce to $O$. The subgroup $D:=\langle\tau\rangle$ is $\Qp$-rational with $D\smallsetminus\{O\}\subset E(\Qp)$, so for $T\in D\smallsetminus\{O\}$ we have $x(T)\in\Zp$ and Theorem~C gives $\phi_p(x(T))=\bar\phi(x(T))^p$. Both sides are units (Step~1 of the proof of Theorem~B), so $\qp(\phi_p(x(T)))=p\,\qp(\bar\phi(x(T)))=0$. The reduction $D\smallsetminus\{O\}\to\Ebar(\Fp)\smallsetminus\{O\}$ is injective ($D\cap\widehat E=0$) between sets of cardinality $p-1$, hence bijective, so $x(T)\bmod p$ runs over all non-zero classes, and Lemma~\ref{lem:local} gives $\lambda\equiv0$.
\end{proof}

\begin{corE}
If $E$ is split at $p$, then $c\ge2$ in every model.
\end{corE}

\begin{proof}
Lemma~\ref{lem:depth} gives $e\ge2$, Corollary~D gives $\lambda=0$, and Theorem~B gives $\qp(\alpha/\beta)=0$.
\end{proof}

The converse of Corollary~D is a known criterion for local torsion, here read as a Fermat quotient.

\begin{prop6}
Let $p\ge5$ be a prime of good ordinary reduction with $\#\Ebar(\Fp)=p$. Then for every $\bar x\in\Ebar(\Fp)\smallsetminus\{O\}$,
\[
\lambda(\bar x)=0\iff E(\Qp)[p]\ne0 .
\]
\end{prop6}

Both sides depend only on $E\bmod p^2$. The equivalence of the right-hand side with ``$E\bmod p^2$ is the canonical lift of $\Ebar$'' is due to Gross \cite{Gro90} (case $n=1$ of the tameness criterion, quoted as \cite[Thm.~4.1]{LR} under the hypotheses $j\ne0,1728$ and $E[p]$ irreducible), and is proved without hypotheses on $j$ in \cite[Lemma~3.2]{DW} and \cite[\S2.1.1, Lemma~4.1]{OY}. The direction $\Leftarrow$ is Corollary~D. For the direction $\Rightarrow$ we give in Appendix~\ref{app:prop6} a proof in terms of $\lambda$ itself, valid for all $p\ge5$: $\lambda$ is an affine function on the deformations of $\Ebar$ over $\Z/p^2$ with a non-zero slope, so it vanishes on exactly one of them, which by Corollary~D is the split one.

\begin{thmS}[the non-split case]
Suppose $\#\Ebar(\Fp)=p$ and $p\nmid n$, and let $P'\in E(\Q)$ have reduction of exact order $p$ and satisfy \textup{(A2)}. If $E(\Qp)[p]=0$, then $\vp\bigl(h_p(pP')\bigr)=1$ exactly, hence $\vp(h_p(P'))=-1$.
\end{thmS}

\begin{proof}
Choose the model with $p\mid a_1$. Since $E$ is not split, $e=1$ by Lemma~\ref{lem:depth}, so $b\ge2$ by Lemma~\ref{lem:decomp}. By Theorem~B, $\qp(\alpha/\beta)=-\lambda$, and $\lambda\ne0$ by Proposition~6, so $c=1$. From $h_p(pP')/2=\Lambda_\sigma+A$ with $b\ge2>1=c$ we get $\vp(h_p(pP'))=1$. Finally $h_p(pP')=p^2h_p(P')$.
\end{proof}

\section{The Main Theorem}\label{sec:main}

\begin{thmx}[Main Theorem]\label{thm:main}
Suppose $\#\Ebar(\Fp)=p$, $p\nmid\prod_\ell c_\ell$ and $\iota$. Let $M_\Lam$ be the height matrix on a basis of $\Lam$.
\begin{enumerate}[label=\textup{(\roman*)}]
\item Every entry of $M_\Lam$ lies in $\Zp$ if and only if $E(\Qp)[p]\ne0$.
\item If $r=1$, then $\vp(\Reg_p)\ge0$ if and only if $E(\Qp)[p]\ne0$, and in the non-split case $\vp(\Reg_p)=-1$.
\item If $E(\Qp)[p]\ne0$, then $\vp(\Reg_p)\ge\max(r-2,0)$.
\end{enumerate}
Without $\iota$, Theorem~\ref{thm:A} already gives $\vp(\Reg_p)\ge r$.
\end{thmx}

\begin{proof}
As $h_p$ is a quadratic form and $p\nmid n$, $\vp(h_p(P))=\vp(h_p(nP))$, and $nP$ satisfies (A2) for every $P\in E(\Q)$.

\emph{($\Leftarrow$ of (i)), $P\in\Lam\smallsetminus\Lam_0$.} The reduction of $P$, hence of $P':=nP$, has exact order $p$, and $P'':=pP'$ satisfies (A1) and (A2). As $E$ is split, Lemma~\ref{lem:depth} gives $e\ge2$, hence $b\ge2$ by Lemma~\ref{lem:decomp} and $c\ge2$ by Corollary~E; so $\vp(h_p(P''))\ge2$, and $h_p(P'')=p^2h_p(P')$ gives $\vp(h_p(P))\ge0$.

\emph{($\Leftarrow$ of (i)), $P\in\Lam_0$.} Then $nP$ satisfies (A1) and (A2), and $\vp(h_p(P))\ge1$ by Lemma~\ref{lem:unit}.

So $\vp(h_p(P))\ge0$ for every $P\in\Lam$; polarising, $h_p(P,R)\in\Zp$ as $p$ is odd. Hence every entry of $M_\Lam$ is integral and $\vp(\Reg_p)\ge0$; with Theorem~\ref{thm:A} this is (iii).

\emph{($\Rightarrow$ of (i)).} Suppose $E(\Qp)[p]=0$. By $\iota$ there is $P\in\Lam$ whose reduction has order $p$; Theorem~S applied to $nP$ gives $\vp(h_p(P))=-1$. If all entries of $M_\Lam$ were integral, so would be the value of the quadratic form at every point of $\Lam$, a contradiction.

\emph{(ii).} If $r=1$ and $Q$ generates $\Lam$, then $Q\notin\Lam_0$ by $\iota$ and $\Reg_p=h_p(Q)$; the two cases above give $\vp(\Reg_p)\ge0$ if $E$ is split and $\vp(\Reg_p)=-1$ otherwise.
\end{proof}

\begin{remark}
The proof needs $\vp(h_p)\ge0$ on $\Lam$, not $\vp(h_p)\ge2$ on $\Lam_0$; for a point of $\Lam\smallsetminus\Lam_0$ it suffices that its reduction have order $p$. That a non-integral entry has valuation \emph{exactly} $-1$, and not merely $\ge-1$, is what Theorem~R below uses.
\end{remark}

\section{The height matrix modulo \texorpdfstring{$p$}{p}}\label{sec:modp}

Throughout this section $\#\Ebar(\Fp)=p$, $p\nmid n$ and $\iota$ hold. Fix a generator $\bar x_0$ of $\Ebar(\Fp)$, and let $\chi\colon\Lam\to\Fp$ be the discrete logarithm of the reduction with respect to it; $\chi$ is onto and $\ker\chi=\Lam_0$.

\subsection{Two scaling properties of $\lambda$}

\begin{thmH}[scaling]
Let $Q\in E(\Zp)$ have non-zero reduction and let $m$ be prime to $p$. Then $\qp\bigl(\phi_p(x(mQ))\bigr)=m^2\,\qp\bigl(\phi_p(x(Q))\bigr)$, that is, $\lambda(m\bar x)=m^2\lambda(\bar x)$.
\end{thmH}

\begin{proof}
The composition rule $\psi_{ab}(R)=\psi_b(R)^{a^2}\psi_a(bR)$ at $(a,b)=(p,m)$ and $(m,p)$, together with $\phi_k(x(R))=x(kR)\psi_k(R)^2$ and the fact that $\phi_m$ is monic of degree $m^2$ over $\Z[a_1,\dots,a_6]$, evaluated at $X:=x(pQ)$ of valuation $-2e\le-2$, gives the identity
\begin{equation}\label{eq:scale}
\phi_p(x(mQ))\,\psi_m(Q)^{2p^2}=\phi_p(x(Q))^{m^2}(1+\delta),\qquad \delta\in p^2\Zp .
\end{equation}
The three division-polynomial values are units, $\psi_m(Q)^{2p^2}$ is a $p$-th power and $1+\delta\in1+p^2\Zp$, so applying $\qp$ gives the claim.
\end{proof}

\begin{corHp}
$\lambda$ vanishes at one non-zero class of $\Ebar(\Fp)$ if and only if it vanishes at all of them; and $\lambda=\lambda_1\chi^2$ on $\Ebar(\Fp)\smallsetminus\{O\}$ with $\lambda_1:=\lambda(\bar x_0)$.
\end{corHp}

\subsection{Bilinear algebra}
Let $K$ be a field of characteristic $\ne2$, $B$ a symmetric bilinear form on a $K$-vector space $V$, $\chi$ a non-zero linear form, $v\in V$ with $\chi(v)=1$, and $W:=\ker\chi$.

\begin{lemma}\label{lem:R0}
\textup{(R0)} If $B(x,x)=0\iff\chi(x)=0$, then $c:=B(v,v)\ne0$ and $B(x,y)=c\,\chi(x)\chi(y)$.

\textup{(R0$_s$)} If $\chi(x)=0\Rightarrow B(x,x)=0$, then
\[
B(x,y)=B(x,v)\chi(y)+\chi(x)B(y,v)-B(v,v)\chi(x)\chi(y);
\]
in particular $\operatorname{rank}B\le2$.
\end{lemma}

\begin{proof}
In both cases $B|_{W\times W}=0$: for $w,w'\in W$, $B(w+w',w+w')=0$ leaves $2B(w,w')=0$. In (R0), moreover, $W\perp v$: if $a:=B(w,v)\ne0$ for some $w\in W$, then $s:=-B(v,v)/(2a)$ makes $x:=sw+v$ satisfy $\chi(x)=1$ and $B(x,x)=0$, a contradiction; and $c\ne0$ since $\chi(v)\ne0$. Writing $x=(x-\chi(x)v)+\chi(x)v$ and similarly for $y$ gives both formulas. In (R0$_s$) the right-hand side involves only $\chi$ and $B(\cdot,v)$, so $\ker\chi\cap\ker B(\cdot,v)$ lies in the radical.
\end{proof}

\subsection{The non-split case}

\begin{thmR}
Suppose in addition $E(\Qp)[p]=0$. Then there is $c\in\Fp^\times$ with
\[
p\,h_p(P,R)\in\Zp\quad\text{and}\quad p\,h_p(P,R)\equiv c\,\chi(P)\chi(R)\pmod p\qquad(P,R\in\Lam).
\]
Hence $\vp(h_p(P,R))=-1$ if and only if $\chi(P)\chi(R)\ne0$, and $\vp(h_p(P,R))\ge1$ when $\chi(P)=\chi(R)=0$. In particular $p\,M_\Lam\bmod p=c\,\chi\chi^t$ has rank exactly one.
\end{thmR}

\begin{proof}
Put $\mathcal Q(P):=p\,h_p(P)$ and $B(P,R):=p\,h_p(P,R)$. For $P\in\Lam_0$, Lemma~\ref{lem:unit} gives $\vp(\mathcal Q(P))\ge2$; for $P\notin\Lam_0$, Theorem~S gives $\vp(\mathcal Q(P))=0$. So $B$ is $\Zp$-valued, and since $p\Lam\subseteq\Lam_0$ it descends to $V:=\Lam/p\Lam$, where the zero set of $x\mapsto\bar B(x,x)$ is exactly the hyperplane $\ker\bar\chi$. Lemma~\ref{lem:R0}(R0) applies. The valuation statements follow by reading the congruence, and Lemma~\ref{lem:unit} polarised for $P,R\in\Lam_0$.
\end{proof}

\begin{propRpp}
In Theorem~R, $c=2\lambda_1$.
\end{propRpp}

\begin{proof}
Let $P'\in\Lam\smallsetminus\Lam_0$ be multiplied by $n$, so that (A2) holds, and $P'':=pP'$; work in a model with $p\mid a_1$. By Lemma~\ref{lem:decomp}, $h_p(P'')=2\Lambda_\sigma+2A$ with $\vp(\Lambda_\sigma)\ge2$, and by Theorem~B and \eqref{eq:logq}, $A\equiv p\lambda\pmod{p^2}$. Hence $h_p(P'')\equiv2p\lambda\pmod{p^2}$ and $p\,h_p(P')\equiv2\lambda\pmod p$. By Corollary~H$'$, $\lambda=\lambda_1\chi^2$, and $c$ is the value at $\chi=1$.
\end{proof}

So in the non-split case the whole matrix $p\,M_\Lam$ modulo $p$ is determined by the single element $\lambda_1\in\Fp$, a Fermat quotient of $\phi_p$ at one point; no rational point enters.

\begin{corRp}
In Theorem~R choose a basis of $\Lam$ with $\chi(e_1)=1$ and $e_2,\dots,e_r\in\Lam_0$. Then $p\,M_\Lam=DA'D$ with $D=\mathrm{diag}(1,p,\dots,p)$ and $A'$ integral with $A'_{11}\in\Zp^\times$; with $T$ the Schur complement of $A'_{11}$ in $A'$,
\[
\vp(\Reg_p)=(r-2)+\vp(\det T).
\]
For $r=1$, $T$ is empty and $\vp(\Reg_p)=-1$.
\end{corRp}

\begin{proof}
By Theorem~R, $p\,h_p(e_1)$ is a unit and $p\,h_p(e_1,e_j)\in p\Zp$ for $j\ge2$; for $i,j\ge2$, Lemma~\ref{lem:unit} polarised gives $p\,h_p(e_i,e_j)\in p^2\Zp$. Hence $p\,M_\Lam=DA'D$ with $A'$ integral, $\det(p\,M_\Lam)=p^{2(r-1)}\det A'$ and $\det(p\,M_\Lam)=p^r\Reg_p$, so $\Reg_p=p^{r-2}\det A'=p^{r-2}A'_{11}\det T$.
\end{proof}

\subsection{The split case}

\begin{thmRs}
Suppose in addition $E(\Qp)[p]\ne0$, and let $\bar M:=M_\Lam\bmod p$, a symmetric bilinear form on $V=\Lam/p\Lam$. Then $\bar M$ vanishes on $\ker\bar\chi\times\ker\bar\chi$, so for any $v$ with $\bar\chi(v)=1$
\[
\bar M(x,y)=\bar M(x,v)\bar\chi(y)+\bar\chi(x)\bar M(y,v)-\bar M(v,v)\bar\chi(x)\bar\chi(y),
\]
and $\rho_1:=\operatorname{rank}\bar M\le2$. In a basis with $\chi(e_1)\ne0$ and $e_2,\dots,e_r\in\Lam_0$, $\bar M=\left(\begin{smallmatrix}a&b^t\\ b&0\end{smallmatrix}\right)$ with $a=h_p(e_1)\bmod p$, $b_j=h_p(e_1,e_j)\bmod p$, and $\rho_1=2\iff b\ne0$.
\end{thmRs}

\begin{proof}
$M_\Lam$ is integral by Theorem~\ref{thm:main}(i). For $P,R\in\Lam_0$, Lemma~\ref{lem:unit} puts $h_p(P),h_p(R),h_p(P+R)$ in $2p\Zp$, so $h_p(P,R)\in p\Zp$; apply Lemma~\ref{lem:R0}(R0$_s$). The leading $2\times2$ minor on $(e_1,e_j)$ is $-b_j^2$.
\end{proof}

\begin{corRps}
In Theorem~R$_s$, $\vp(\Reg_p)\ge r-\rho_1$. If $\rho_1=2$, then $\vp(\Reg_p)=(r-2)+\vp(\det T_s)$ for an integral $(r-2)\times(r-2)$ Schur complement $T_s$. For $r=2$: $\vp(\Reg_p)=0\iff\rho_1=2$.
\end{corRps}

\begin{proof}
Write $M:=M_\Lam$ in the basis of Theorem~R$_s$, so $M_{ij}\in p\Zp$ for $i,j\ge2$. If $\rho_1=2$, a change of basis inside $\Lam_0\otimes\Zp$ (which does not change $\vp(\det M)$) brings $b$ to $(b_1,0,\dots,0)$ with $b_1$ a unit. The block $S=\left(\begin{smallmatrix}M_{11}&b_1\\ b_1&pC\end{smallmatrix}\right)$ has $\det S\equiv-b_1^2\not\equiv0$, so $S\in GL_2(\Zp)$ and its Schur complement in $M$ is $p\,T_s$ with $T_s$ integral; hence $\det M=\det S\cdot p^{r-2}\det T_s$. If $\rho_1\le1$, then $p\mid M_{1j}$ for $j\ge2$ as well, so rows $2,\dots,r$ are divisible by $p$ and $\vp(\Reg_p)\ge r-1$; if moreover $a=0$, then $\vp(\Reg_p)\ge r$.
\end{proof}

Since $\rho_1\le2$, Corollary~R$'_s$ contains Theorem~\ref{thm:main}(iii) and says when it is attained. The constant $a$ is again a Fermat quotient, one digit deeper. In the split case $\lambda=0$ by Corollary~D, so for $u=\phi_p(P')$ we have $u^{p-1}\in1+p^2\Zp$, and the \emph{second Fermat quotient}
\[
q^{(2)}(u):=\frac{u^{p-1}-1}{p^2}\bmod p\qquad(u\in\Zp^\times,\ \qp(u)=0)
\]
is defined; it is a homomorphism on that subgroup.

\begin{propRpps}
In Theorem~R$_s$ let $e_1\in\Lam$ with $\chi(e_1)\ne0$, $P':=ne_1$, and $\lambda^{(2)}:=\tfrac12q^{(2)}(\phi_p(P'))\in\Fp$. Then $a=h_p(e_1)\bmod p=2\lambda^{(2)}/n^2$. In particular, for $r=1$, $\vp(\Reg_p)=0\iff\lambda^{(2)}\ne0$.
\end{propRpps}

\begin{proof}
Work in a model with $p\mid a_1$, put $P'':=pP'$ and $e:=\vp(d(P''))\ge2$ (Lemma~\ref{lem:depth}). In Step~2 of the proof of Theorem~B, every term of $\omega^2-\phi^3$ other than $a_1\phi\omega\psi$ has valuation $\ge2e\ge4$, and that one has valuation $\ge e+1\ge3$; hence $\omega^2\equiv\phi^3\pmod{p^3}$. Corollary~D gives $\qp(\phi)=\qp(\omega)=0$; writing $\phi^{p-1}=1+p^2A_2$ and $\omega^{p-1}=1+p^2B_2$ we get $2B_2\equiv3A_2\pmod p$, so $q^{(2)}(\phi)=2\lambda^{(2)}$ and $q^{(2)}(\omega)=3\lambda^{(2)}$. The factor $d'^{p^2}$ of Step~1 has $q^{(2)}=0$, so $q^{(2)}(\alpha/\beta)=-\lambda^{(2)}$. For a unit $u$ with $\qp(u)=0$ one has $\log_pu\equiv-p^2q^{(2)}(u)\pmod{p^3}$, so $A\equiv p^2\lambda^{(2)}\pmod{p^3}$, while $\vp(\Lambda_\sigma)\ge3$ by Lemma~\ref{lem:decomp}. Hence $h_p(P'')\equiv2p^2\lambda^{(2)}\pmod{p^3}$, and $h_p(P'')=p^2n^2h_p(e_1)$.
\end{proof}

The analogy with Proposition~R$''$ is exact, with one difference: $\lambda$ depends only on $E\bmod p^2$ and the class $\bar x$, whereas $\lambda^{(2)}$ depends on $E\bmod p^3$ and on the point modulo $p^2$, because $\phi_p'$ is divisible by $p$ but not by $p^2$. There is therefore no law for the split value of $\vp(\Reg_p)$ in terms of the reduction alone; its finer structure is studied in \cite{Companion}.

\subsection{The case $\#\Ebar(\F_5)=10$}

\begin{propRplus}
Let $p\ge5$ have good ordinary reduction with $p\mid N:=\#\Ebar(\Fp)$ and $p\nmid n$, let $\chi\colon\Lam\to\Ebar(\Fp)[p]\cong\Z/p$ be the $p$-part of the reduction, and assume $\chi$ onto. Then Theorem~S, Theorem~R, Proposition~R$''$ and Corollary~R$'$ hold as stated. Moreover $p\,\|\,N$, and the only case not covered by $\#\Ebar(\Fp)=p$ is $N=2p$, which forces $p=5$.
\end{propRplus}

\begin{proof}
If $\chi(P)=0$, the order of the reduction of $P$ divides $m:=N/p$, so $mnP$ satisfies (A1) and (A2) and $h_p(P)=h_p(mnP)/(mn)^2\in2p\Zp$. If $\chi(P)\ne0$, then $mnP$ has reduction of exact order $p$ and Theorem~S applies to it. The bilinear algebra is unchanged. Finally $N\le p+1+2\sqrt p<p^2$ gives $N\in\{p,2p\}$, and $N=2p$ forces $(p-1)^2\le4p$, hence $p=5$.
\end{proof}

\section{The CM case and rank-two CM curves}\label{sec:CM}

\begin{thmF}[{\cite[Prop.~2.1(3)]{DW}}]
Suppose $E/\Q$ has complex multiplication and $p\ge5$ is a prime of good ordinary reduction with $a_p=1$. Then $E(\Qp)[p]\ne0$.
\end{thmF}

This is due to David and Weston, who prove it for $p>3$ together with the converse. For completeness we sketch a self-contained argument. The prime $p$ splits in the CM field $k$, so the CM order acts on $E_{/\Qp}$; by Deuring's lifting theorem Frobenius lifts to an endomorphism $\pi$ with $\pi+\bar\pi=1$ and $\pi\bar\pi=p$. The isogeny $\varrho:=1-\pi=\bar\pi$ has degree $p$ and is separable on the special fibre, as $(1-F)^*\omega=\omega$; so $\ker\varrho$ is étale of order $p$ over $\Zp$, Frobenius acts on it as $\pi=1$, and its points are therefore defined over $\Qp$.

\begin{corFp}
If $E$ has CM, $p\ge5$ is anomalous of good ordinary reduction, $p\nmid n$ and $\iota$ holds, then $\vp(\Reg_p)\ge\max(r-2,0)$. A CM curve never shows the deficit.
\end{corFp}

Recent work of Banwait \cite{Ban} proves, for rank-two CM curves and outside an explicit set of primes, that the normalised second Taylor coefficient of the Mazur--Tate--Teitelbaum $p$-adic $L$-function is a unit if and only if the cyclotomic regulator is a unit and $\mathrm{Sha}(E/\Q)[p^\infty]=0$; the excluded set contains the anomalous primes. Our results describe the regulator at those primes.

\begin{corG}
Let $E/\Q$ have CM and rank two, and let $p\ge5$ be an anomalous prime of good ordinary reduction with $p\nmid n$ and $\iota$. Then the height matrix is $p$-integral, $\vp(\Reg_p)\ge0$, and
\[
\vp(\Reg_p)=0\iff \operatorname{rank}_{\Fp}(M_\Lam\bmod p)=2\iff h_p(e_1,e_2)\not\equiv0\pmod p
\]
for any basis with $\chi(e_2)=0$.
\end{corG}

\begin{proof}
By Theorem~F, $E$ is split at $p$; apply Theorem~\ref{thm:main}(i) and Corollary~R$'_s$ with $r=2$.
\end{proof}

For curves with CM by $\Z[i]$ at most one prime is anomalous \cite[Lemma~2.2]{Ban}, so the corollary concerns very few primes for each curve; but it shows that at these primes the regulator is integral and that whether it is a unit is decided by a single entry of the height matrix modulo $p$. To extend the criterion of \cite{Ban} itself to such a prime one would also have to account for the Euler factor $(1-\alpha^{-1})^2$ of the $p$-adic $L$-function, which has valuation $2$ at an anomalous prime; we do not pursue this here.

\section{Related work}\label{sec:related}

\emph{The baseline.} The negative valuations at anomalous primes appear in the tables of \cite[\S\S4.2--4.6]{MST} (for instance $37a$ at $p=53$, where the anomaly is noted); the factor $n_{\mathrm{lcm}}=\operatorname{lcm}(\#\Ebar(\Fp),\operatorname{lcm}_\ell c_\ell)$ in the normalisation of \cite{Har} costs $2\vp(n_{\mathrm{lcm}})$ digits of precision at an anomalous prime; and \cite[\S2.2]{Ban} uses non-anomalousness to pass from $E^0(\Q)$ to $E(\Q)$. Integrality of the pairing at non-anomalous primes not dividing the Tamagawa product is used in \cite{BCS}. We did not find Theorem~\ref{thm:A} stated as a proposition, and we do not claim it.

\emph{The local criterion.} Proposition~6 is the case $n=1$ of Gross's criterion \cite{Gro90} as quoted in \cite[Thm.~4.1]{LR}, and is contained in \cite[Lemma~3.2]{DW} and \cite[Lemma~4.1]{OY}. The same condition appears in cryptography: the attacks of Semaev, Smart and Satoh--Araki on anomalous curves \cite{Sem98,Sma99,SA98} fail exactly for lifts congruent to the canonical lift modulo $p^2$ \cite{Vol2}, so Corollary~D is the vanishing of a Fermat quotient on the canonical lift. Silverman's and Voloch's elliptic Wieferich conditions \cite{Sil88,Vol} concern $N_pP$ modulo $p^2$ for a rational point and are different. We did not find the criterion in the form $\qp(\phi_p(x_0))=0$.

\emph{$p$-adic heights.} Wuthrich studied the variation of $p$-adic heights in families and formulated conjectures on the valuation of the $p$-adic regulator \cite{Wut04}; the quantity $e=\vp(t(pP'))$ of Section~\ref{sec:criterion} also appears in \cite{Wut}. We did not find the splitting hypothesis $E(\Qp)[p]\ne0$ in connection with the valuation of $\Reg_p$, nor a statement that the height pairing is $p$-integral exactly under it.

\appendix

\section{Proof of Proposition 6}\label{app:prop6}

The direction $\Leftarrow$ is Corollary~D. For $\Rightarrow$ we show that $\lambda$, viewed as a function on the lifts of $\Ebar$ to $\Z/p^2$, is affine with non-zero slope; since exactly one lift is split (Serre--Tate), and $\lambda$ vanishes there by Corollary~D, it vanishes nowhere else. The arithmetic steps are formalised in Lean (Appendix~\ref{app:lean}); the analytic steps use the Tate parametrisation \cite{ATAEC,Dok}.

\emph{(1) An Euler identity.} The polynomial $\phi_p$ is weighted homogeneous of weight $2p^2$ for the weights $(4,6,2)$ on $(a_4,a_6,x)$ in short Weierstrass form $y^2=x^3+a_4x+a_6$. Writing $\phi_p\equiv f(X)$ with $X=x^p$ of weight $2p$ and reducing $4a_4\partial_{a_4}\phi_p+6a_6\partial_{a_6}\phi_p+2x\partial_x\phi_p=2p^2\phi_p$ modulo $p$ gives $4a_4\,\partial_{a_4}f+6a_6\,\partial_{a_6}f=0$ in $\Fp[a_4,a_6][X]$. Hence $a_4\mid\partial_{a_6}f$; put $\tilde u:=(\partial_{a_6}f)/a_4$.

\emph{(2) One slope.} For the deformation $E_s\colon y^2=x^3+(a_4+sp)x+a_6$, Taylor expansion of $\phi_p$ in its coefficients gives $\lambda_s=\lambda_0-s\,B_4/2$, and likewise a slope in the $a_6$ direction. The invariance of $\lambda$ under $(x,y)\mapsto(u^2x,u^3y)$ (under which $\phi_p$ is multiplied by a $p$-th power), applied with $u=1+pt$, relates the two slopes by the identity of (1), so both are multiples of one scalar $\rho(\bar x)=-\tilde u(\bar x)/(4f(\bar x))$; the denominator is non-zero by Vélu's formulas. The family $E_s$ exhausts the deformations of $\Ebar$ over $\Z/p^2$, and exactly one of them is split (Serre--Tate theory). An affine function with non-zero slope has exactly one zero, so it remains to show $\rho\ne0$.

\emph{(3) A dichotomy.} Let $g$ be the polynomial with $\psi_p\equiv g(x^p)$ modulo $p$, of degree $(p-1)/2$, whose roots are the $x$-coordinates of $\ker V\smallsetminus\{O\}$. The resultant $\mathrm{Res}_X(\tilde u,g)$ is weighted homogeneous of weight $12m$ with $m:=(p-1)(p-2)(p+3)/12$. Along the family $E_t\colon y^2=x^3-3x+(2+t)$ over $\bar\F_p((t))$, which has split multiplicative reduction with $v(q)=v(\Delta)=1$, the Tate parametrisation gives $\ker V=\langle q\rangle\bmod q^p$; the roots $\beta_j$ of $g$ satisfy $v(\beta_j-1)=j$ for $1\le j\le(p-1)/2$, and the Newton polygon of $f$ gives $v(f(\beta_j))=j(p-j)$. Since $\tilde u=-\tfrac13\,\frac{d}{dt}f$, one obtains $\operatorname{ord}_\Delta\mathrm{Res}_X(\tilde u,g)\ge\sum_j(j(p-j)-1)=m$, and weighted homogeneity forces $\mathrm{Res}_X(\tilde u,g)=c\,\Delta^m$ with a constant $c\in\Fp$. Thus either $\rho\ne0$ for every ordinary curve over $\bar\F_p$ and every point of $\ker V\smallsetminus\{O\}$, or $\rho=0$ for all of them.

\emph{(4) Which side.} The Frobenius twist of $E_t$ is the member at parameter $t^p$, so $d/dt$ acts only through the Tate parameter and $v(d\beta_j/dt)=j-1$. Among the roots of $f$, the two nearest to $\beta_j$ contribute $1/(1-w)+1/(1-w^{-1})=1$ to $f'(\beta_j)/f(\beta_j)$, which therefore has valuation exactly $-j$ with leading coefficient $p-2j\equiv-2j$. The chain rule gives
\[
12\rho_j=\frac{d}{dt}\log f(\beta_j)-\frac{f'(\beta_j)}{f(\beta_j)}\frac{d\beta_j}{dt}=\frac{-j^2+2j^2}{t}+O(1),
\]
so the residues $-j^2$ and $-2j^2$ do not cancel, $v(\rho_j)=-1$, $\operatorname{ord}_\Delta\mathrm{Res}_X(\tilde u,g)=m$ exactly, and $c\ne0$. Hence $\rho\ne0$ for every ordinary curve, which proves Proposition~6. \qed

\section{Formalisation}\label{app:lean}

The Lean development accompanying this paper consists of eleven files, each checking from \texttt{import Mathlib} alone (Lean~4 \texttt{v4.33.1}, Mathlib commit \texttt{0df444a3}); \texttt{sorry} and \texttt{native\_decide} do not occur and every \texttt{\#print axioms} line is a subset of \{\texttt{propext}, \texttt{Classical.choice}, \texttt{Quot.sound}\}. The files and the axiom logs are in the repository \cite{Repo}. Lean checks that the stated propositions follow from their hypotheses; the geometric inputs listed in the introduction are hypotheses of the Lean statements, not conclusions.

\begin{center}
\small
\begin{longtable}{p{0.32\linewidth}p{0.48\linewidth}p{0.1\linewidth}}
\toprule
Statement & Lean declarations & File\\
\midrule
Theorem~\ref{thm:A}, Lemma~\ref{lem:decomp} & \lean{norm_det_div_sq_le}, \lean{dvd_pow_sub_one}, \lean{fermatQuot_additive} & 1\\
Lemma~\ref{lem:local} & \lean{eval_sub_eval_dvd_sq} & 2\\
Theorem B & \lean{fermatQuot_sub_iff_of_sq_eq_cube}, \lean{sq_sub_cube_dvd_of_dvd_a1}, \lean{sq_sub_cube_dvd_of_sq_dvd_psi}, \lean{dvd_three_mul_sub_two_mul} & 2, 8\\
Corollaries D, E & \lean{sq_dvd_pow_sub_one_of_isPow}, \lean{sq_dvd_alpha_sub_beta_of_isPow} & 3\\
Theorem H, Corollary H$'$ & \lean{fermatQuot_scale}, \lean{fermatQuot_scale_zero_iff} & 4\\
Theorem~\ref{thm:main} (linear algebra) & \lean{regulator_valuation_of_split_anomalous}, \lean{regulator_valuation_of_split_anomalous_iff_split} & 5\\
Theorem~\ref{thm:main} ($\Rightarrow$ of (i)) & \lean{quadratic_value_mem_of_entries_mem} & 8\\
Appendix A (2) & \lean{slope_of_rho}, \lean{fermatQuot_deform}, \lean{affine_zero_unique}, \lean{lambda_eq_zero_iff_split} & 5\\
Appendix A (3) & \lean{weight_split}, \lean{twelve_dvd_weight}, \lean{sum_pole_orders}, \lean{cusp_sum_eq_m}, \lean{eq_C_mul_of_dvd_of_natDegree_eq} & 6\\
Appendix A (4) & \lean{first_term_residue}, \lean{second_term_leading_ne_zero}, \lean{residues_do_not_cancel} & 7\\
Theorem S & \lean{exists_shift_dvd_a1}, \lean{add_val_exactly_one}, \lean{not_dvd_of_val_one} & 8\\
Theorem R, Corollary R$'$ & \lean{bilin_rank_one_of_zero_iff}, \lean{det_conj_diagonal_one_p}, \lean{reg_eq_of_det_eq} & 9\\
Theorem R$_s$, Corollary R$'_s$, Proposition R$''_s$ & \lean{bilin_eq_of_vanishing_on_hyperplane}, \lean{gram_det_three_eq_zero}, \lean{det_two_dvd_iff}, \lean{pow_dvd_det_of_offdiag_block}, \lean{fermatQuot2_add}, \lean{split_const_eq} & 10\\
Proposition R$^+$ (counting) & \lean{card_cases}, \lean{two_p_forces_five} & 22\\
\bottomrule
\end{longtable}
\end{center}
File $k$ is \texttt{PadicHeightValuation}$k$\texttt{.lean} (file 1 has no suffix).

\section*{Acknowledgements and use of AI}
The author carried out this research with the assistance of Claude Code (Anthropic). The constructions, proofs, computations and the Lean~4 formalisation were developed with the AI system under the author's direction; independent sessions of the system were used to re-check the Lean development and the computations. The author reviewed the arguments, checked the Lean development and its axiom logs, and takes full responsibility for the content of this paper.

\begin{thebibliography}{MTT86}

\bibitem[ATAEC]{ATAEC} J.~H. Silverman, \emph{Advanced Topics in the Arithmetic of Elliptic Curves}, Grad. Texts in Math.\ 151, Springer, 1994.
\bibitem[Ban]{Ban} B.~S. Banwait, \emph{Second derivatives of $p$-adic $L$-functions and the Shafarevich--Tate group of rank-two CM elliptic curves}, arXiv:2609.08431 (2026).
\bibitem[BCS]{BCS} J.~S. Balakrishnan, M.~\c{C}iperiani, W.~Stein, \emph{$p$-adic heights of Heegner points and $\Lambda$-adic regulators}, Math. Comp.\ 84 (2015), 923--954.
\bibitem[BMS]{BMS} J.~S. Balakrishnan, J.~S. M\"uller, W.~A. Stein, \emph{A $p$-adic analogue of the conjecture of Birch and Swinnerton-Dyer for modular abelian varieties}, Math. Comp.\ 85 (2016), 983--1016.
\bibitem[Comp]{Companion} Y.~Hosomura, \emph{The $p$-adic height at split anomalous primes: the digits of the height and an archimedean threshold}, companion paper (2026).
\bibitem[Dok]{Dok} T.~Dokchitser, V.~Dokchitser, \emph{Local invariants of isogenous elliptic curves}, Trans. Amer. Math. Soc.\ 367 (2015), 4339--4358.
\bibitem[Draft]{Draft} Kiichi (Y.~Hosomura), \emph{Integrality of the $p$-adic height pairing at an anomalous prime is equivalent to a point of order $p$ in $E(\Qp)$}, draft v15, computo ergo sum (2026), Zenodo, doi:10.5281/zenodo.23296340.
\bibitem[DW]{DW} C.~David, T.~Weston, \emph{Local torsion on elliptic curves and the deformation theory of Galois representations}, Math. Res. Lett.\ 15 (2008), 599--611.
\bibitem[Gro90]{Gro90} B.~H. Gross, \emph{A tameness criterion for Galois representations associated to modular forms (mod $p$)}, Duke Math. J.\ 61 (1990), 445--517.
\bibitem[Har]{Har} D.~Harvey, \emph{Efficient computation of $p$-adic heights}, LMS J. Comput. Math.\ 11 (2008), 40--59.
\bibitem[LR]{LR} \'A.~Lozano-Robledo, \emph{Division fields of elliptic curves with minimal ramification}, Rev. Mat. Iberoam.\ 31 (2015), 1311--1332.
\bibitem[Maz72]{Maz72} B.~Mazur, \emph{Rational points of abelian varieties with values in towers of number fields}, Invent. Math.\ 18 (1972), 183--266.
\bibitem[MST]{MST} B.~Mazur, W.~Stein, J.~Tate, \emph{Computation of $p$-adic heights and log convergence}, Doc. Math.\ Extra Vol.\ (2006), 577--614.
\bibitem[MT91]{MT91} B.~Mazur, J.~Tate, \emph{The $p$-adic sigma function}, Duke Math. J.\ 62 (1991), 663--688.
\bibitem[MTT]{MTT} B.~Mazur, J.~Tate, J.~Teitelbaum, \emph{On $p$-adic analogues of the conjectures of Birch and Swinnerton-Dyer}, Invent. Math.\ 84 (1986), 1--48.
\bibitem[OY]{OY} Y.~Ozeki, M.~Yoshida, \emph{Fields of definition of $p$-torsion points of elliptic curves and their ramification}, arXiv:2609.06554 (2026).
\bibitem[Repo]{Repo} Y.~Hosomura, \emph{Lean development and data for this paper}, GitHub repository and Zenodo record (to be inserted).
\bibitem[SA98]{SA98} T.~Satoh, K.~Araki, \emph{Fermat quotients and the polynomial time discrete log algorithm for anomalous elliptic curves}, Comment. Math. Univ. St. Pauli 47 (1998), 81--92.
\bibitem[Sch85]{Sch85} P.~Schneider, \emph{$p$-adic height pairings II}, Invent. Math.\ 79 (1985), 329--374.
\bibitem[Sem98]{Sem98} I.~A. Semaev, \emph{Evaluation of discrete logarithms in a group of $p$-torsion points of an elliptic curve in characteristic $p$}, Math. Comp.\ 67 (1998), 353--356.
\bibitem[Sil88]{Sil88} J.~H. Silverman, \emph{Wieferich's criterion and the abc-conjecture}, J. Number Theory 30 (1988), 226--237.
\bibitem[Sma99]{Sma99} N.~P. Smart, \emph{The discrete logarithm problem on elliptic curves of trace one}, J. Cryptology 12 (1999), 193--196.
\bibitem[SW]{SW} W.~Stein, C.~Wuthrich, \emph{Algorithms for the arithmetic of elliptic curves using Iwasawa theory}, Math. Comp.\ 82 (2013), 1757--1792.
\bibitem[Vol]{Vol} J.~F. Voloch, \emph{Elliptic Wieferich primes}, J. Number Theory 81 (2000), 205--209.
\bibitem[Vol2]{Vol2} J.~F. Voloch, \emph{Relating the Smart--Satoh--Araki and Semaev approaches to the discrete logarithm problem on anomalous elliptic curves}, note.
\bibitem[Wut]{Wut} C.~Wuthrich, \emph{On $p$-adic elliptic logarithms and $p$-adic approximation lattices}, 2006.
\bibitem[Wut04]{Wut04} C.~Wuthrich, \emph{On $p$-adic heights in families of elliptic curves}, J. London Math. Soc.\ (2) 70 (2004), 23--40.

\end{thebibliography}
\end{document}
